% ID: SEL-D0071-OA-E03
% ESTADO: validada. S071.
\textbf{Parte I.} Por inclusión y exclusión, $P(A\cap B)=0{,}30+0{,}40-0{,}65=0{,}05$.

\textbf{a)} Son compatibles porque la intersección tiene probabilidad positiva.

\textbf{b)} Son dependientes porque $0{,}05\ne P(A)P(B)=0{,}12$.

\textbf{c)} La parte de $A$ fuera de $B$ tiene probabilidad $0{,}30-0{,}05=0{,}25$. Al condicionar a $B^c$ resulta $P(A\mid B^c)=0{,}25/0{,}60=\boxed{5/12}$.

\textbf{Parte II. a)} Los nueve valores suman $64{,}8$, por lo que $\bar x=7{,}2$. La desviación es conocida y el error estándar es $0{,}9/\sqrt9=0{,}3$. Con $z_{0{,}985}\approx2{,}17009$,
\[
7{,}2\pm z_{0{,}985}\cdot0{,}3
\approx\boxed{[6{,}549,\ 7{,}851]}.
\]
\textbf{b)} La semiamplitud impresa es $(7{,}494-6{,}906)/2=0{,}294$. Igualamos a $z_{0{,}975}\,0{,}9/\sqrt n$ y despejamos:
\[
n=\left(\frac{z_{0{,}975}\,0{,}9}{0{,}294}\right)^2\approx36.
\]
El tamaño es $\boxed{36}$. Para ese tamaño, el margen redondea a $0{,}294$, como en el intervalo dado.
